%%═══════════════════════════════════════════════════════════════════
%% Lecture 19 — 1D Hydrodynamics: Finite Differences on Staggered Grid
%% 77315 — Computational Physics
%%═══════════════════════════════════════════════════════════════════

\section{Lecture 19 — 1D Hydrodynamics (Finite Differences)}\label{sec:lec19}
\hrule\vspace{1.5em}

%%──────────────────────────────────────────────────────────────────
\subsection*{Overview}
We discretise the Lagrangian 1D Euler equations on a
\emph{staggered grid}: thermodynamic quantities ($\rho$, $\epsilon$, $p$)
live at cell centres while velocity and acceleration live at cell
\emph{edges} (nodes).  Artificial viscosity smooths shocks over
3–5 cells.  The Courant condition limits the time step.

%%──────────────────────────────────────────────────────────────────
\subsection{Staggered Grid Layout}\label{subsec:lec19:stagger}

\begin{center}
\begin{tabular}{ccccc}
  $\cdot$ & $|\,\rho_i,\epsilon_i,p_i\,|$ & $\cdot$
          & $|\,\rho_{i+1},\epsilon_{i+1},p_{i+1}\,|$ & $\cdot$ \\
  $u_{i-1/2}$ & & $u_{i+1/2}$ & & $u_{i+3/2}$
\end{tabular}
\end{center}

\begin{itemize}
  \item \textbf{Cells} $i = 1,\ldots,N$: hold $\rho_i$, $\epsilon_i$, $p_i$, mass $m_i$.
  \item \textbf{Nodes} $i+\tfrac{1}{2}$: hold velocity $u_{i+1/2}$ and
        coordinate $x_{i+1/2}$.
\end{itemize}

%%──────────────────────────────────────────────────────────────────
\subsection{Update Equations per Time Step}\label{subsec:lec19:update}

The eight update equations (using superscript $n$ for old,
$n+1$ for new):

\begin{enumerate}
  \item \textbf{Equation of motion} (momentum, at nodes):
        \begin{equation}
          u_{i+1/2}^{n+1} = u_{i+1/2}^n
          - \frac{\Delta t}{m_i + m_{i+1}}\cdot 2A_{i+1/2}
            \left(p_{i+1}^n - p_i^n + q_{i+1}^n - q_i^n\right),
          \label{eq:lec19:eom}
        \end{equation}
        where $q_i$ is the artificial viscosity pressure and $A_{i+1/2}$
        is the cell-interface area.

  \item \textbf{Coordinate update} (advance node positions):
        \begin{equation}
          x_{i+1/2}^{n+1} = x_{i+1/2}^n + \Delta t\, u_{i+1/2}^{n+1}.
        \end{equation}

  \item \textbf{Volume update} (cell volume from node positions).

  \item \textbf{Density update} (from mass conservation):
        $\rho_i^{n+1} = m_i / V_i^{n+1}$.

  \item \textbf{Velocity divergence} (dilatation):
        $(\nabla\cdot\mathbf{u})_i = (u_{i+1/2}^{n+1}-u_{i-1/2}^{n+1})/\Delta x_i$.

  \item \textbf{Artificial viscosity update} (see below).

  \item \textbf{Energy update} (1st law):
        \begin{equation}
          \epsilon_i^{n+1} = \epsilon_i^n
          - \Delta t\,(p_i^n + q_i^n)\,(\nabla\cdot\mathbf{u})_i.
          \label{eq:lec19:energy}
        \end{equation}

  \item \textbf{Pressure update} via EOS: $p_i^{n+1} = (\gamma-1)\rho_i^{n+1}\epsilon_i^{n+1}$.
\end{enumerate}

%%──────────────────────────────────────────────────────────────────
\subsection{Artificial Viscosity}\label{subsec:lec19:artvis}

Physical shocks are discontinuities that would alias on any finite
grid.  Introduce a \emph{pseudo-pressure} $q$ that acts only during
compression ($\nabla\cdot\mathbf{u} < 0$):
\begin{equation}
  q_i = \begin{cases}
    C_q\,\rho_i\,(\Delta u_i)^2 & \text{if } \Delta u_i < 0 \text{ (compression)},\\
    0 & \text{otherwise},
  \end{cases}
  \label{eq:lec19:artvis}
\end{equation}
where $\Delta u_i = u_{i+1/2} - u_{i-1/2}$ and $C_q \sim 2$
(von Neumann–Richtmyer viscosity).  This smears the shock over
$\sim 3$–5 cells, enabling stable integration through it.

The Hugoniot jump conditions (Rankine–Hugoniot) are automatically
satisfied to first order once the artificial viscosity is included.

%%──────────────────────────────────────────────────────────────────
\subsection{Courant Condition}\label{subsec:lec19:courant}

The time step must satisfy
\begin{equation}
  \Delta t \le \frac{C_\text{CFL}\,\Delta x_\text{min}}{c_s + |u|_\text{max}},
  \quad C_\text{CFL} \approx 0.2\text{--}0.5.
  \label{eq:lec19:courant}
\end{equation}
Typical codes use $C_\text{CFL} = 0.2$ to maintain stability with
the artificial viscosity term.

\nt{The staggered grid is crucial: centering density/pressure and
velocity on different grids avoids the checkerboard instability that
plagues collocated schemes for compressible flow.}
